\section{The original central curve carries the canonical leading field and order}

Let \(i:C\hookrightarrow Y\) be a genuine closed immersion, \(x\in C\),
\(S=\mathcal O_{Y,i(x)}\), and \(A=\mathcal O_{C,x}\).
The original component equation \(\pi\in S\) satisfies
\(\ker(i_x^\#)=(\pi)\). It is the equation of the chosen reduced component,
not the product defining a whole crossing central fibre.

\begin{definition}[The actual original central stalk quotient]
\label{def:central-curve-original-stalk-quotient}
\uses{def:fc-closed-immersion-principal-stalk-quotient,thm:fc-closed-immersion-stalk-order-length}
\lean{EnvelopeProperty.CentralCurveLeadingOrders.stalkQuotientEquiv,EnvelopeProperty.CentralCurveLeadingOrders.stalkQuotientEquiv_mk}
\leanok
There is the canonical ring isomorphism
\(e:S/(\pi)\xrightarrow{\sim}A\).
For every original \(s\in S\), it satisfies
\(e(\bar s)=i_x^\#(s)\).
No field, residue, or order comparison is supplied.
\end{definition}
\begin{proof}
The actual stalk map of the genuine closed immersion is surjective.
Identify its actual kernel with the supplied original principal component
ideal, and apply the first isomorphism theorem to that same map.
The quotient-class formula is the literal first-isomorphism formula.
This is the stalk-kernel form of the existing affine-equation quotient
construction; it requires neither a new choice of a stalk map nor a
codimension-one assertion about the ambient closed point.
\end{proof}

Assume now \(C\) is integral and \(S\) is a domain. Retain the original
parameters \(z,\pi\) with \(\pi\) prime, \(\pi\nmid z\), and
\(\mathfrak m_S=(\pi,z)\). Put \(R=S/(\pi)\), \(T=S_{(\pi)}\).

\begin{definition}[Canonical leading field of the same original curve]
\label{def:central-curve-original-leading-field}
\uses{def:central-curve-original-stalk-quotient,lem:gauss-surface-canonical-vertical-leading-field}
\lean{EnvelopeProperty.CentralCurveLeadingOrders.residualFractionEquiv,EnvelopeProperty.CentralCurveLeadingOrders.residualFractionEquiv_coefficient,EnvelopeProperty.CentralCurveLeadingOrders.verticalResidueCurveFieldEquiv,EnvelopeProperty.CentralCurveLeadingOrders.verticalResidueCurveFieldEquiv_coefficient}
\leanok
The actual quotient isomorphism induces
\(F:\operatorname{Frac}R\xrightarrow{\sim}C.\mathrm{functionField}\).
Composing with the existing canonical
\(E_v:\kappa(T)\xrightarrow{\sim}\operatorname{Frac}R\) gives
\(E_C:\kappa(T)\xrightarrow{\sim}C.\mathrm{functionField}\).
For every \(s\in S\), both maps retain the actual coefficient diagram:
\[
 F(\bar s)=\iota_A(i_x^\#s),\qquad
 E_C(\operatorname{res}_T(s))=\iota_A(i_x^\#s).
\]
The field \(\kappa(T)\) here is a function field, not the closed-point field
\(\kappa_C(x)\).
\end{definition}
\begin{proof}
The intrinsic stalk-to-generic-stalk map already makes the same original
\(C.\mathrm{functionField}\) the fraction field of \(A\).
Extend the internally constructed whole ring isomorphism \(e\) by the
canonical fraction-ring universal property. Its coefficient law and the
preceding quotient-class formula give the first displayed identity.
Compose with the existing canonical vertical residue equivalence.
Its already proved every-coefficient law gives the second identity.
No independent fraction-field map or coefficient-matching certificate is
introduced.
\end{proof}

\begin{lemma}[Actual fraction orders under a whole ring isomorphism]
\label{lem:central-curve-fraction-order-ring-equiv}
\uses{thm:fc-original-principal-ring-iso-order}
\lean{EnvelopeProperty.CentralCurveLeadingOrders.ordFrac_ringEquiv}
\leanok
Let \(R_0,A_0\) be Noetherian domains of Krull dimension at most one,
with actual fraction fields \(K_0,L_0\).
For a genuine ring isomorphism \(e_0:R_0\simeq A_0\), let \(F_0:K_0\simeq L_0\)
be its canonical induced fraction-field isomorphism. For every \(f\in K_0\),
including zero,
\[
 \operatorname{ordFrac}_{R_0}(f)
 =\operatorname{ordFrac}_{A_0}(F_0(f)).
\]
\end{lemma}
\begin{proof}
The already proved semilinear principal-quotient isomorphism preserves
the actual module lengths over the respective original rings.
Thus \(\operatorname{ord}_{R_0}(a)=\operatorname{ord}_{A_0}(e_0(a))\).
For a nonzero \(a\), the native coefficient formula for
\(\operatorname{ordFrac}\) expresses both sides by the same canonical
embedding of that length; zero follows from preservation of zero.
Write any actual \(f\) as a quotient of original coefficients with nonzero
denominator. The canonical coefficient commutation and preservation of
division extend the equality to every fraction.
This gives exact normalized orders, without an unknown ramification factor.
\end{proof}

Assume for the remaining statements that \(S\) is Noetherian, \(C\) is
locally Noetherian, \(\dim C\le1\), and \(x\ne\eta_C\).

\begin{theorem}[The same curve's full fractional leading order]
\label{thm:central-curve-original-leading-order}
\uses{def:central-curve-original-leading-field,lem:central-curve-fraction-order-ring-equiv,lem:gauss-surface-original-residual-vertical-dvrs,thm:fc-curve-proper-support-points}
\lean{EnvelopeProperty.CentralCurveLeadingOrders.residualOrder_eq_curveOrdHom,EnvelopeProperty.CentralCurveLeadingOrders.leadingResidue_order_eq_curveOrd}
\leanok
For every \(r\in\operatorname{Frac}R\), the old residual order is precisely
the original curve order homomorphism:
\[
 \operatorname{residualOrder}(r)
 = C.\operatorname{ordHom}_x(F(r)).
\]
For every actual unit \(u\in T^\times\), its actual leading image is nonzero
and
\[
 \operatorname{residualOrder}(E_v(\operatorname{res}_T u))
 =\exp\!\left(C.\operatorname{ord}_x(E_C(\operatorname{res}_T u))\right).
\]
\end{theorem}
\begin{proof}
The original curve dimension and nongeneric point give actual codimension
one at \(x\). The old original parameter construction makes \(R\) a DVR,
hence a Noetherian domain of dimension one.
Apply the preceding exact fraction-order transport to the constructed
quotient isomorphism \(e:R\simeq A\).
The definition of the original curve's order homomorphism is the native
fraction order of its same stalk \(A\), proving the first identity.
The residue of an actual unit of \(T\) is a unit of its residue field.
It remains nonzero under both actual field isomorphisms.
Apply the original nonzero-function comparison between
\(C.\operatorname{ordHom}_x\) and \(\exp(C.\operatorname{ord}_x)\).

The first statement deliberately uses \(\operatorname{ordHom}\) at zero:
the integer order at zero has a junk value and cannot give this exponential
identity there. Leading images of vertical units can have poles at \(x\);
their integer curve order can therefore be negative.
\end{proof}

\begin{theorem}[Original eventual orders with the actual curve coefficient]
\label{thm:central-curve-original-eventual-order}
\uses{thm:central-curve-original-leading-order,thm:gauss-surface-original-fraction-eventual-orders}
\lean{EnvelopeProperty.CentralCurveLeadingOrders.exists_eventual_fieldOrder_curve}
\leanok
For every nonzero original \(f\in\operatorname{Frac}S\), internally produce
\(a\in\mathbf Z\) and an actual \(u\in T^\times\) such that
\[
 f=\iota_T(u)\pi^a,\qquad \operatorname{verticalOrder}(f)=\exp(a),
\]
and, for all sufficiently large original \(N>0\),
\[
 \operatorname{weightedOrder}_N(f)
 =\exp\!\left(aN+C.\operatorname{ord}_x(E_C(\operatorname{res}_T u))\right).
\]
All fraction embeddings and weighted DVRs are the unchanged old ones.
\end{theorem}
\begin{proof}
Apply the existing all-fraction eventual theorem once. It internally gives
the original \(a,b,u\), the factorization, the vertical order, the actual
leading residue order \(\exp(b)\), and the eventual weighted order
\(\exp(aN+b)\). The preceding true original-curve identity identifies that
same leading residue order with the exponential of the same curve's native
integer order. Injectivity of the integer exponential gives the equality
with \(b\). Substitute it in the already proved eventual tail.
This consumes the old theorem without assuming a leading-field/order
comparison and without rebuilding its valuation family.
\end{proof}

